\begin{myexercise}
    Stellen Sie sich vor, Ihre Schule nutzt eine Lern-Plattform namens "Schul-Net", über die Sie Nachrichten an Lehrkräfte schreiben und Hausaufgaben hochladen könnt.
    In den folgenden drei Szenarien wurde jeweils eines der drei Sicherheitsziele verletzt. 
    Für jedes Szenario:
    \begin{itemize}
        \item Analysieren Sie den Vorfall: Welches der drei Sicherheitsziele (Vertraulichkeit, Integrität, Authentifikation) wurde primär verletzt?
        \item Beschreiben Sie, warum die (in diesem Szenario) noch funktionierenden zwei Ziele das Problem nicht verhindern konnten.
    \end{itemize}
    
    \textbf{Szenario A}: Herr Roth schickt über ``Schul-Net'' an jeden seiner Schüler ein individuelles, privates
    Feedback zu ihren Leistungen im letzten Quartal. Die Noten werden dabei nicht über das System, sondern nur in dieser privaten Nachricht mitgeteilt.\\
    Am nächsten Tag wird in der Pause auf dem Schulhof über die genauen Noten und auch über einige sehr spezifische, peinliche Rechtschreibfehler von Mitschüler Ben diskutiert, die so nur in Herrn Roths privatem Feedback an Ben stehen konnten. Ben ist sich sicher, dass niemand sein Passwort hat und er sich immer ausgeloggt hat.

    \textbf{Szenario B}: Für das Geschichtsreferat lädt Schülerin Anna ihre fertige Präsentationsdatei (eine
    .pptx-Datei) in den Abgabe-Ordner von ``Schul-Net'' hoch. Sie hat die Datei auf Viren geprüft, bevor sie sie
    hochgeladen hat.\\
    Als der Lehrer, Herr Roth, die Datei eine Stunde später für die Präsentation herunterlädt und öffnet, ist die Präsentation intakt. Auf der Titelfolie ist Annas Name, und alle Folien sind da. Während Anna jedoch präsentiert, erscheint auf Folie 5 plötzlich ein animiertes Werbe-Banner für ein Online-Casino, das Anna dort definitiv nicht eingefügt hat. Sie ist blamiert.

    \textbf{Szenarien C}: Die Schule führt die Wahl zum Schülersprecher über ``Schul-Net'' durch. Das System ist so
    aufgebaut, dass jede Stimme anonym abgegeben wird und die Stimmen korrekt gezählt werden. Nach der Wahl gewinnt ein eher unbekannter Kandidat mit einer überwältigenden Mehrheit. Das Ergebnis ist so überraschend, dass die Schulleitung eine Umfrage in den Klassen macht. Es stellt sich heraus, dass viele Schüler berichten, sie hätten gar nicht an der Wahl teilgenommen – im System ist aber vermerkt, dass von ihren Accounts eine Stimme abgegeben wurde.

    \begin{myanswer}
        \textbf{Szenario A}: Verletztes Ziel: Vertraulichkeit.\\
        Begründung: Private Informationen (Bens Note und Feedback) sind an die Öffentlichkeit gelangt. Obwohl die
        Nachricht integer (unverändert) war und die Authentifikation (Ben und Herr Roth) stimmte, war der Kanal
        selbst nicht vertraulich. Jemand konnte die Kommunikation ``abhören'' (z.B. durch einen
        Man-in-the-Middle-Angriff oder eine Lücke im Server, die Einblick in fremde Postfächer erlaubte).
Ω
        \textbf{Szenario B}: Verletztes Ziel: Integrität.\\
        Begründung: Die Datei wurde verändert, nachdem Anna sie erstellt hatte. Die Daten sind nicht mehr in ihrem
        korrekten, ursprünglichen Zustand. Obwohl die Vertraulichkeit gewahrt wurde (niemand hat den Inhalt
        gestohlen, nur manipuliert) und die Authentifikation funktionierte (Anna hat es hochgeladen, Herr Roth hat
        es heruntergeladen), war die Datei selbst nicht mehr ``unversehrt''.

        \textbf{Szenario C}: Verletztes Ziel: Authentifikation.\\
        Begründung: Dies ist das kniffligste Szenario. Das Wahlergebnis fühlt sich ``falsch'' an (wie ein
        Integritätsproblem). Aber die Stimmen, nachdem sie abgegeben wurden, wurden ja korrekt gezählt (Integrität)
        und blieben geheim (Vertraulichkeit). Das Kernproblem ist, dass das System nicht sicherstellen konnte, dass
        die Person, die auf ``Wählen'' klickte, wirklich der/die berechtigte Schüler/in war. Ein Angreifer konnte die
        Identität (Authentifikation) anderer Schüler vortäuschen oder übernehmen und in deren Namen abstimmen
    \end{myanswer}
\end{myexercise}